\chapter{2014年天津卷（文）}

\section{单选题}

\begin{problem}
\(\i\) 是虚数单位，复数 \(\frac{7 + \i}{3 + 4\i} =\)（\quad）

\choices
    {\(1 - \i\)}
    {\(-1 + \i\)}
    {\(\frac{17}{25} + \frac{31}{25}\i\)}
    {\(-\frac{17}{7} + \frac{25}{7}\i\)}
\end{problem}

\answer{A}

\begin{solution}
分子、分母同乘以分母的共轭复数，得
\[
\frac{7+\i}{3+4\i}=\frac{(7+\i)(3-4\i)}{(3+4\i)(3-4\i)}=\frac{21-28\i+3\i-4\i^{2}}{25}=\frac{25-25\i}{25}=1-\i
\]
因此答案为 A.
\end{solution}

\begin{problem}
设变量 \(x\)，\(y\) 满足约束条件 \(\begin{cases}
x + y - 2 \geqslant 0,\\
x - y - 2 \leqslant 0,\\
y \geqslant 1,
\end{cases}\) 则目标函数 \(z = x + 2y\) 的最小值为（\quad）

\choices
    {\(2\)}
    {\(3\)}
    {\(4\)}
    {\(5\)}
\end{problem}

\answer{B}

\begin{solution}
约束条件等价于 \(x\geqslant2-y\)，\(x\leqslant y+2\)，\(y\geqslant1\). 对固定的 \(y\)，\(z=x+2y\) 随 \(x\) 增大而增大，所以取允许的最小值 \(x=2-y\). 此时 \(z=2-y+2y=2+y\geqslant3\). 当 \(y=1\)，\(x=1\) 时各约束均成立，且 \(z=3\)，故最小值为 \(3\)，答案为 B.
\end{solution}

\begin{problem}
已知命题 \( p: \forall x > 0 \)，总有 \( (x + 1)\e^x > 1 \)，则 \( \neg p \) 为（\quad）

\choices
    {\(\exists x_0 \leqslant 0\)，使得 \((x_0 + 1)\e^{x_0} \leqslant 1\)}
    {\(\exists x_0 > 0\)，使得 \((x_0 + 1)\e^{x_0} \leqslant 1\)}
    {\(\forall x > 0\)，总有 \((x + 1)\e^{x} \leqslant 1\)}
    {\(\forall x \leqslant 0\)，总有 \((x + 1)\e^{x} \leqslant 1\)}
\end{problem}

\answer{B}

\begin{solution}
命题 \(p\) 是在条件 \(x>0\) 下的全称命题，全称命题“对所有 \(x>0\)，不等式成立”的否定是“存在一个 \(x_0>0\)，使不等式不成立”，即 \((x_0+1)\e^{x_0}\leqslant1\). 因此答案为 B.
\end{solution}

\begin{problem}
设 \(a = \log_2\pi\)，\(b = \log_{\frac{1}{2}}\pi\)，\(c = \pi^{-2}\)，则（\quad）

\choices
    {\(a > b > c\)}
    {\(b > a > c\)}
    {\(a > c > b\)}
    {\(c > b > a\)}
\end{problem}

\answer{C}

\begin{solution}
因为 \(\pi>2\)，所以 \(a=\log_2\pi>1\). 又因为底数 \(\frac12\) 在 \(0\) 与 \(1\) 之间且真数 \(\pi>1\)，所以 \(b=\log_{1/2}\pi<0\). 另外 \(c=\pi^{-2}\) 满足 \(0<c<1\). 因而 \(a>c>b\)，答案为 C.
\end{solution}

\begin{problem}
设 \(\{a_n\}\) 是首项为 \(a_1\)，公差为 \(-1\) 的等差数列，\(S_n\) 为其前 \(n\) 项和. 若 \(S_1\)，\(S_2\)，\(S_4\) 成等比数列，则 \(a_1 =\)（\quad）

\choices
    {\(2\)}
    {\(-2\)}
    {\(\frac{1}{2}\)}
    {\(-\frac{1}{2}\)}
\end{problem}

\answer{D}

\begin{solution}
由公差为 \(-1\)，得 \(S_1=a_1\)，\(S_2=2a_1-1\)，\(S_4=4a_1-6\). 三项成等比数列，所以
\[
(2a_1-1)^2=a_1(4a_1-6)
\]
展开并化简得 \(4a_1^2-4a_1+1=4a_1^2-6a_1\)，即 \(2a_1+1=0\). 因此 \(a_1=-\frac12\)，答案为 D.
\end{solution}

\begin{problem}
已知双曲线 \(\frac{x^2}{a^2} -\frac{y^2}{b^2} = 1\)（\(a > 0\)，\(b > 0\)）的一条渐近线平行于直线 \(l: y = 2x + 10\)，双曲线的一个焦点在直线 \(l\) 上，则双曲线的方程为（\quad）

\choices
    {\(\frac{x^2}{5} -\frac{y^2}{20} = 1\)}
    {\(\frac{x^2}{20} -\frac{y^2}{5} = 1\)}
    {\(\frac{3x^2}{25} -\frac{3y^2}{100} = 1\)}
    {\(\frac{3x^2}{100} -\frac{3y^2}{25} = 1\)}
\end{problem}

\answer{A}

\begin{solution}
双曲线的渐近线斜率为 \(\pm\frac ba\). 由一条渐近线平行于 \(l\)，得 \(\frac ba=2\)，即 \(b=2a\). 设半焦距为 \(c\)，则 \(c^2=a^2+b^2=5a^2\). 双曲线的焦点为 \((\pm c,0)\)，在直线 \(l\) 上的焦点满足 \(0=2x+10\)，故该焦点为 \((-5,0)\)，从而 \(c=5\). 因此 \(a^2=5\)，\(b^2=20\)，双曲线方程为 \(\frac{x^2}{5}-\frac{y^2}{20}=1\)，答案为 A.
\end{solution}

\begin{problem}
如图，\(\triangle ABC\) 是圆的内接三角形，\(\angle BAC\) 的平分线交圆于点 \(D\)，交 \(BC\) 于 \(E\)，过点 \(B\) 的圆的切线与 \(AD\) 的延长线交于点 \(F\).

\bitmapfigure[max width=5.8cm]{img/2014/tianjin_liberal/q7_fig1.png}

在上述条件下，给出下列四个结论：

\begin{circlelist}
\item \(BD\) 平分 \(\angle CBF\)；
\item \(FB^{2} = FD \cdot FA\)；
\item \(AE \cdot CE = BE \cdot DE\)；
\item \(AF \cdot BD = AB \cdot BF\).
\end{circlelist}

则所有正确结论的序号是（\quad）

\choices
    {\circled{1}\circled{2}}
    {\circled{3}\circled{4}}
    {\circled{1}\circled{2}\circled{3}}
    {\circled{1}\circled{2}\circled{4}}
\end{problem}

\answer{D}

\begin{solution}
设 \(\angle BAD=\angle DAC=\alpha\).

由圆周角定理，\(\angle DBC=\angle DAC=\alpha\). 由切线与弦所夹的角等于弦所对的圆周角，\(\angle FBD=\angle BAD=\alpha\). 因而 \(BD\) 平分 \(\angle CBF\)，结论 \(1\) 正确.

又因为 \(F\)，\(A\)，\(D\) 共线，且由同样的弦切角关系有 \(\angle FBD=\angle FAB\)，同时 \(\angle BFD=\angle BFA\)，所以 \(\triangle FBD\sim\triangle FAB\). 对应边给出
\[
\frac{FB}{FA}=\frac{FD}{FB}=\frac{BD}{AB}
\]
于是 \(FB^2=FD\cdot FA\)，且 \(AF\cdot BD=AB\cdot BF\)，结论 \(2\)，\(4\) 正确.

对于结论 \(3\)，点 \(E\) 是两弦 \(AD\)，\(BC\) 的交点，由相交弦定理实际有
\[
AE\cdot DE=BE\cdot CE
\]
并非题给的等式. 例如取等边三角形，设外接圆半径为 \(R\)，则 \(E\) 为 \(BC\) 的中点，\(AE=\frac{3R}{2}\)，\(DE=\frac R2\)，\(BE=CE=\frac{\sqrt3R}{2}\). 此时 \(AE\cdot CE=\frac{3\sqrt3R^2}{4}\)，而 \(BE\cdot DE=\frac{\sqrt3R^2}{4}\)，二者不等，故结论 \(3\) 错误.

所以所有正确结论的序号为 \(1\)，\(2\)，\(4\)，答案为 D.
\end{solution}

\begin{problem}
已知函数 \( f(x) = \sqrt{3}\sin \omega x + \cos \omega x \)（\(\omega > 0\)），\(x \in \R\). 在曲线 \( y = f(x) \) 与直线 \( y = 1 \) 的交点中，若相邻交点距离的最小值为 \( \frac{\pi}{3} \)，则 \( f(x) \) 的最小正周期为（\quad）

\choices
    {\(\frac{\pi}{2}\)}
    {\(\frac{2\pi}{3}\)}
    {\(\pi\)}
    {\(2\pi\)}
\end{problem}

\answer{C}

\begin{solution}
利用辅助角公式，\(f(x)=2\sin\left(\omega x+\frac\pi6\right)\). 交点满足
\[
\sin\left(\omega x+\frac\pi6\right)=\frac12
\]
因此 \(\omega x=2k\pi\) 或 \(\omega x=\frac{2\pi}{3}+2k\pi\)，其中 \(k\in\Z\). 相邻交点的横坐标间距交替为 \(\frac{2\pi}{3\omega}\) 和 \(\frac{4\pi}{3\omega}\)，故最小值为 \(\frac{2\pi}{3\omega}=\frac\pi3\)，解得 \(\omega=2\). 所以最小正周期为 \(T=\frac{2\pi}{\omega}=\pi\)，答案为 C.
\end{solution}

\section{填空题}

\begin{problem}
某大学为了解在校本科生对参加某项社会实践活动的意向，拟采用分层抽样的方法，从该校四个年级的本科生中抽取一个容量为 \(300\) 的样本进行调查. 已知该校一年级，二年级，三年级，四年级的本科生人数之比为 \(4:5:5:6\)，则应从一年级本科生中抽取\fillinblank{} 名学生.
\end{problem}

\answer{60}

\begin{solution}
四个年级人数的总份数为 \(4+5+5+6=20\). 分层抽样中各层抽取人数之比等于各层人数之比，所以一年级应抽取
\[
300\times\frac4{20}=60
\]
名学生.
\end{solution}

\begin{problem}
已知一个几何体的三视图如图所示（单位：\(\mathrm{m}\)），则该几何体的体积为\fillinblank{}\(\mathrm{m}^{3}\).

\bitmapfigure[max width=7.2cm]{img/2014/tianjin_liberal/q10_fig1.png}
\end{problem}

\answer{\(\frac{20\pi}{3}\)}

\begin{solution}
由三视图可知，该几何体由圆柱和圆锥组成：圆柱的底面直径为 \(2 \mathrm{~m}\)，高为 \(4 \mathrm{~m}\)，圆锥的底面直径为 \(4 \mathrm{~m}\)，高为 \(2 \mathrm{~m}\). 因而圆柱、圆锥的底面半径分别为 \(1 \mathrm{~m}\)，\(2 \mathrm{~m}\)，体积为
\[
V=\pi\cdot(1 \mathrm{~m})^2\cdot4 \mathrm{~m}+\frac13\pi\cdot(2 \mathrm{~m})^2\cdot2 \mathrm{~m}=\frac{20\pi}{3}\ \mathrm{m}^{3}
\]
\end{solution}

\begin{problem}
阅读如图的框图，运行相应的程序，输出 \(S\) 的值为\fillinblank{}.

\bitmapfigure[max width=4.8cm]{img/2014/tianjin_liberal/q11_fig1.png}
\end{problem}

\answer{\(-4\)}

\begin{solution}
初始时 \(S=0\)，\(n=3\). 第一次循环后，\(S=0+(-2)^3=-8\)，\(n=2\)，此时 \(n\leqslant1\) 不成立，继续循环. 第二次循环后，\(S=-8+(-2)^2=-4\)，\(n=1\)，此时 \(n\leqslant1\) 成立，程序输出 \(S=-4\).
\end{solution}

\begin{problem}
函数 \( f(x) = \lg x^{2} \) 的单调递减区间是\fillinblank{}.
\end{problem}

\answer{\((-\infty,0)\)}

\begin{solution}
函数定义域为 \(x\ne0\). 当 \(x<0\) 时，\(x^2\) 随 \(x\) 增大而减小，而 \(\lg t\) 在 \((0,+\infty)\) 上单调递增，所以 \(\lg x^2\) 在 \((-\infty,0)\) 上单调递减. 当 \(x>0\) 时，\(x^2\) 与 \(x\) 同增，函数单调递增. 因此单调递减区间为 \((-\infty,0)\).
\end{solution}

\begin{problem}
已知菱形 \(ABCD\) 的边长为 \(2\)，\(\angle BAD = 120^{\circ}\)，点 \(E\)，\(F\) 分别在边 \(BC\)，\(DC\) 上，\(BC = 3BE\)，\(DC = \lambda DF\). 若 \(\overrightarrow{AE} \cdot \overrightarrow{AF} = 1\)，则 \(\lambda\) 的值为\fillinblank{}.
\end{problem}

\answer{2}

\begin{solution}
设 \(\bs u=\overrightarrow{AB}\)，\(\bs v=\overrightarrow{AD}\). 则 \(|\bs u|=|\bs v|=2\)，且
\[
\bs u\cdot\bs u=\bs v\cdot\bs v=4
\]
又有 \(\bs u\cdot\bs v=4\cos120^{\circ}=-2\).
由 \(BC=3BE\)，\(DC=\lambda DF\)，得
\[
\overrightarrow{AE}=\bs u+\frac13\bs v
\]
因为 \(F\) 在 \(DC\) 上且 \(\overrightarrow{DF}=\frac1\lambda\overrightarrow{DC}=\frac1\lambda\bs u\)，所以
\[
\overrightarrow{AF}=\overrightarrow{AD}+\overrightarrow{DF}=\bs v+\frac1\lambda\bs u
\]
于是
\[
\overrightarrow{AE}\cdot\overrightarrow{AF}
=\left(\bs u+\frac13\bs v\right)\cdot\left(\bs v+\frac1\lambda\bs u\right)
=-2+\frac4\lambda+\frac43-\frac{2}{3\lambda}
=\frac{10}{3\lambda}-\frac23
\]
由题设该数量积等于 \(1\)，所以 \(\frac{10}{3\lambda}-\frac23=1\)，解得 \(\lambda=2\).
\end{solution}

\begin{problem}
已知函数 \( f(x) = \begin{cases}
|x^2 + 5x + 4|, & x \leqslant 0,\\
2|x - 2|, & x > 0
\end{cases} \) 若函数 \( y = f(x) - a|x| \) 恰有 \(4\) 个零点，则实数 \(a\) 的取值范围为\fillinblank{}.
\end{problem}

\answer{\((1,2)\)}

\begin{solution}
零点方程为 \(f(x)=a|x|\). 先看 \(x>0\). 当 \(0<x<2\) 时，方程为 \(4-2x=ax\)，有根 \(x=\frac4{a+2}\)；当 \(x>2\) 时，方程为 \(2x-4=ax\)，有根 \(x=\frac4{2-a}\). 当 \(0<a<2\) 时这两个根分别位于 \((0,2)\) 与 \((2,+\infty)\)；当 \(a=2\) 只有前一个根，当 \(a>2\) 没有这两个正根.

对 \(x\leqslant0\)，令 \(t=-x\geqslant0\)，则方程化为
\[
\left|t^2-5t+4\right|=at
\]
当 \(0\leqslant t\leqslant1\) 或 \(t\geqslant4\) 时，令 \(p(t)=t^2-(5+a)t+4\)，方程为 \(p(t)=0\). 对于 \(a>0\)，有 \(p(0)=4>0\)，\(p(1)=-a<0\)，\(p(4)=-4a<0\)，且当 \(t\to+\infty\) 时 \(p(t)\to+\infty\)，所以 \(p(t)\) 在 \((0,1)\) 和 \((4,+\infty)\) 内各有一个根；由于 \(p\) 是二次函数，这两个根就是该方程的全部根.

当 \(1\leqslant t\leqslant4\) 时，方程为 \(g(t)=t^2-(5-a)t+4=0\). 由于 \(g(1)=a>0\)，\(g(4)=4a>0\)，且当 \(0<a<1\) 时顶点 \(\frac{5-a}{2}\in(2,2.5)\)，并且 \(g\left(\frac{5-a}{2}\right)=4-\frac{(5-a)^2}{4}<0\)，所以此时在 \([1,4]\) 内有两个不同的根. 当 \(a=1\) 时，\(g(t)=(t-2)^2\)，有一个重根；当 \(a>1\) 时，若 \(1<a<5\)，则顶点值非负，若 \(a\geqslant5\)，则 \(g'(t)=2t-(5-a)>0\) 在 \([1,4]\) 上恒成立，结合 \(g(1)>0\)，可知区间 \([1,4]\) 内无根.

因此，当 \(0<a<1\) 时左侧有 \(4\) 个根、右侧有 \(2\) 个根，共 \(6\) 个；当 \(a=1\) 时共有 \(2+1+2=5\) 个；当 \(1<a<2\) 时左侧有 \(2\) 个、右侧有 \(2\) 个，共 \(4\) 个；当 \(a\geqslant2\) 时正半轴只有 \(1\) 个根，左半轴有 \(2\) 个根，共 \(3\) 个. 当 \(a=0\) 时，方程 \(f(x)=0\) 的根为 \(-4\)，\(-1\)，\(2\)，共 \(3\) 个；当 \(a<0\) 时右端 \(a|x|<0\) 而 \(f(x)\geqslant0\)，不可能产生所需的四个零点，故 \(a\) 的取值范围为 \((1,2)\).
\end{solution}

\section{解答题}

\begin{problem}
某校夏令营有 \(3\) 名男同学 \(A\)，\(B\)，\(C\) 和 \(3\) 名女同学 \(X\)，\(Y\)，\(Z\)，其年级情况如下表.

\begin{center}
\begin{tabular}{c|ccc}
\hline
 & 一年级 & 二年级 & 三年级 \\
\hline
男同学 & \(A\) & \(B\) & \(C\) \\
\hline
女同学 & \(X\) & \(Y\) & \(Z\) \\
\hline
\end{tabular}
\end{center}

现从这 \(6\) 名同学中随机选出 \(2\) 人参加知识竞赛（每人被选到的可能性相同）.

\begin{enumerate}
\item 用表中字母列举出所有可能的结果；
\item 设 \(M\) 为事件“选出的 \(2\) 人来自不同年级且恰有 \(1\) 名男同学和 \(1\) 名女同学”，求事件 \(M\) 发生的概率.
\end{enumerate}
\end{problem}

\begin{solution}
\begin{enumerate}
\item 由于不计选取顺序，所有等可能结果为
\[
\begin{gathered}
(A,B)\quad(A,C)\quad(A,X)\quad(A,Y)\quad(A,Z)\quad(B,C)\quad(B,X)\quad(B,Y)\quad(B,Z)\quad\\
(C,X)\quad(C,Y)\quad(C,Z)\quad(X,Y)\quad(X,Z)\quad(Y,Z)\quad
\end{gathered}
\]
共 \(\mathrm C_6^2=15\) 个.
\item 事件 \(M\) 包含的结果为 \((A,Y)\)，\((A,Z)\)，\((B,X)\)，\((B,Z)\)，\((C,X)\)，\((C,Y)\)，共 \(6\) 个. 因此
\[
P(M)=\frac6{15}=\frac25
\]
\end{enumerate}
\end{solution}

\begin{problem}
在 \(\triangle ABC\) 中，内角 \(A\)，\(B\)，\(C\) 所对的边分别为 \(a\)，\(b\)，\(c\)，已知 \(a - c = \frac{\sqrt{6}}{6} b\)，\(\sin B = \sqrt{6} \sin C\).

\begin{enumerate}
\item 求 \(\cos A\) 的值；
\item 求 \( \cos\left(2A-\frac{\pi}{6}\right) \) 的值.
\end{enumerate}
\end{problem}

\begin{solution}
由正弦定理，\(\frac b c=\frac{\sin B}{\sin C}=\sqrt6\)，所以 \(b=\sqrt6c\). 代入已知条件得
\[
a-c=\frac{\sqrt6}{6}\cdot\sqrt6c=c
\]
从而 \(a=2c\). 由余弦定理，
\[
\cos A=\frac{b^2+c^2-a^2}{2bc}=\frac{6c^2+c^2-4c^2}{2\sqrt6c^2}=\frac{\sqrt6}{4}
\]
由于 \(A\) 是三角形内角，\(\sin A=\frac{\sqrt{10}}4\). 于是
\(\cos2A=2\cos^2A-1=-\frac14\)，\(\sin2A=2\sin A\cos A=\frac{\sqrt{15}}4\).
因此
\[
\cos\left(2A-\frac\pi6\right)=\cos2A\cos\frac\pi6+\sin2A\sin\frac\pi6=\frac{\sqrt{15}-\sqrt3}{8}
\]
\end{solution}

\begin{problem}
如图，四棱锥 \(PABCD\) 的底面 \(ABCD\) 是平行四边形，\(BA = BD = \sqrt{2}\)，\(AD = 2\)，\(PA = PD = \sqrt{5}\)，\(E\)，\(F\) 分别是棱 \(AD\)，\(PC\) 的中点.

\begin{enumerate}
\item 证明：\(EF \parallel\) 平面 \(PAB\)；
\item 若二面角 \(P - AD - B\) 为 \(60^{\circ}\)，

\begin{enumerate}
\item 证明：平面 \(PBC \perp\) 平面 \(ABCD\)；
\item 求直线 \(EF\) 与平面 \(PBC\) 所成角的正弦值.
\end{enumerate}
\end{enumerate}

\bitmapfigure[max width=7.3cm]{img/2014/tianjin_liberal/q17_fig1.png}
\end{problem}

\begin{solution}
\begin{enumerate}
\item 设平行四边形两条对角线的交点为 \(H\). 则 \(H\) 是 \(AC\) 的中点，也是 \(BD\) 的中点. 在 \(\triangle ADC\) 中，\(E\)，\(H\) 分别是 \(AD\)，\(AC\) 的中点，所以 \(EH\parallel DC\parallel AB\). 在 \(\triangle PAC\) 中，\(F\)，\(H\) 分别是 \(PC\)，\(AC\) 的中点，所以 \(FH\parallel PA\). 因而平面 \(EFH\parallel\) 平面 \(PAB\)，从而 \(EF\parallel\) 平面 \(PAB\).
\item
\begin{enumerate}
\item 因为 \(BA^2+BD^2=AD^2\)，所以 \(\triangle ABD\) 在 \(B\) 处为直角三角形. 由 \(E\) 是斜边 \(AD\) 的中点，得 \(BE=1\). 又因 \(BA=BD\)，\(BE\) 是等腰三角形 \(ABD\) 底边 \(AD\) 上的中线，故 \(BE\perp AD\). 又 \(PA=PD\)，\(PE\) 是等腰三角形 \(PAD\) 底边 \(AD\) 上的中线，故 \(PE\perp AD\)，并且
\[
PE=\sqrt{PA^2-AE^2}=\sqrt{5-1}=2
\]
二面角 \(P-AD-B\) 的平面角为 \(\angle PEB\)，故 \(\angle PEB=60^{\circ}\). 在 \(\triangle PBE\) 中，
\[
PB^2=PE^2+BE^2-2PE\cdot BE\cos60^{\circ}=4+1-2=3
\]
在 \(\triangle PBD\) 中，\(PB^2+BD^2=3+2=5=PD^2\)，所以 \(PB\perp BD\). 在 \(\triangle PAB\) 中，\(PB^2+BA^2=3+2=5=PA^2\)，所以 \(PB\perp BA\). 因为 \(BA\)，\(BD\) 是平面 \(ABCD\) 内相交直线，故 \(PB\perp\) 平面 \(ABCD\). 又 \(PB\subset\) 平面 \(PBC\)，所以平面 \(PBC\perp\) 平面 \(ABCD\).
\item 取 \(B\) 为原点，分别以 \(BA\)，\(BD\)，\(BP\) 所在直线为坐标轴. 由上可知三轴两两垂直，且
\(A(\sqrt2,0,0)\)，\(D(0,\sqrt2,0)\)，\(P(0,0,\sqrt3)\).
由平行四边形的向量关系 \(\overrightarrow{BC}=\overrightarrow{AD}\)，得 \(C(-\sqrt2,\sqrt2,0)\). 因而
\(E\left(\frac{\sqrt2}{2},\frac{\sqrt2}{2},0\right)\)，\(F\left(-\frac{\sqrt2}{2},\frac{\sqrt2}{2},\frac{\sqrt3}{2}\right)\)
所以直线 \(EF\) 的方向向量可取
\[
\bs d=\overrightarrow{EF}=\bigl(-\sqrt2,0,\frac{\sqrt3}{2}\bigr)
\]
平面 \(PBC\) 内的两个方向向量可取 \(\overrightarrow{PB}=(0,0,-\sqrt3)\)，\(\overrightarrow{BC}=(-\sqrt2,\sqrt2,0)\)，故其法向量可取 \(\bs n=(1,1,0)\). 设直线 \(EF\) 与平面 \(PBC\) 所成角为 \(\theta\)，则
\[
\sin\theta=\frac{|\bs d\cdot\bs n|}{|\bs d|\,|\bs n|}=\frac{\sqrt2}{\frac{\sqrt{11}}2\cdot\sqrt2}=\frac2{\sqrt{11}}=\frac{2\sqrt{11}}{11}
\]
\end{enumerate}
\end{enumerate}
\end{solution}

\begin{problem}
设椭圆 \(\frac{x^2}{a^2} + \frac{y^2}{b^2} = 1\) （\(a > b > 0\)）的左，右焦点分别为 \(F_1\)，\(F_2\)，右顶点为 \(A\)，上顶点为 \(B\)，已知 \(|AB| = \frac{\sqrt{3}}{2} |F_1F_2|\).

\begin{enumerate}
\item 求椭圆的离心率；
\item 设 \(P\) 为椭圆上异于其顶点的一点，以线段 \(PB\) 为直径的圆经过点 \(F_{1}\)，经过点 \(F_{2}\) 的直线 \(l\) 与该圆相切于点 \(M\)，\(|MF_{2}| = 2\sqrt{2}\). 求椭圆的方程.
\end{enumerate}
\end{problem}

\begin{solution}
设椭圆半焦距为 \(c\)，则 \(c^2=a^2-b^2\)，且 \(|AB|=\sqrt{a^2+b^2}\)，\(|F_1F_2|=2c\). 由已知
\[
a^2+b^2=3c^2=3(a^2-b^2)\quad
\]
所以 \(a^2=2b^2\)，\(c^2=b^2\)，从而
\[
e=\frac ca=\frac{\sqrt2}{2}
\]

由 \(b=c\)，椭圆可写为 \(\frac{x^2}{2c^2}+\frac{y^2}{c^2}=1\)，并且 \(B=(0,c)\)，\(F_1=(-c,0)\)，\(F_2=(c,0)\). 设 \(P=(x,y)\). 因为以 \(PB\) 为直径的圆经过 \(F_1\)，所以 \(\angle PF_1B=90^{\circ}\)，即
\[
(x+c,y)\cdot(c,c)=0\quad\text{从而}\quad x+y=-c
\]
联立椭圆方程，令 \(u=\frac xc\)，得
\(\frac{u^2}{2}+(u+1)^2=1\)，\(u(3u+4)=0\).
当 \(u=0\) 时，\(P=(0,-c)\) 是椭圆的顶点，应舍去；故 \(P=\left(-\frac{4c}{3},\frac c3\right)\).

圆心 \(O\) 是 \(PB\) 的中点，所以 \(O=\left(-\frac{2c}{3},\frac{2c}{3}\right)\)，半径满足 \(r=\frac{|PB|}{2}=\frac{\sqrt5c}{3}\). 因此
\[
OF_2^2=\left(\frac{5c}{3}\right)^2+\left(\frac{2c}{3}\right)^2=\frac{29c^2}{9}
\]
直线 \(l\) 从 \(F_2\) 引圆的切线，故 \(OM\perp F_2M\)，由勾股定理
\[
MF_2^2=OF_2^2-OM^2=\frac{29c^2}{9}-\frac{5c^2}{9}=\frac{8c^2}{3}
\]
代入 \(MF_2=2\sqrt2\)，得 \(8=\frac{8c^2}{3}\)，所以 \(c^2=3\). 从而 \(a^2=6\)，\(b^2=3\)，椭圆方程为 \(\frac{x^2}{6}+\frac{y^2}{3}=1\).
\end{solution}

\begin{problem}
已知函数 \( f(x) = x^2 - \frac{2}{3} a x^3 \)（\(a > 0\)），\(x \in \R\).

\begin{enumerate}
\item 求 \( f(x) \) 的单调区间和极值；
\item 若对于任意的 \(x_{1} \in (2, +\infty)\)，都存在 \(x_{2} \in (1, +\infty)\)，使得 \(f(x_{1}) \cdot f(x_{2}) = 1\)，求 \(a\) 的取值范围.
\end{enumerate}
\end{problem}

\begin{solution}
\begin{enumerate}
\item
\[
f'(x)=2x-2ax^2=2x(1-ax)
\]
当 \(x<0\) 时 \(f'(x)<0\)；当 \(0<x<\frac1a\) 时 \(f'(x)>0\)；当 \(x>\frac1a\) 时 \(f'(x)<0\). 所以 \(f(x)\) 在 \((-\infty,0)\) 和 \(\left(\frac1a,+\infty\right)\) 上单调递减，在 \(\left(0,\frac1a\right)\) 上单调递增. 又
\(f(0)=0\)，\(f\left(\frac1a\right)=\frac1{a^2}-\frac{2}{3a^2}=\frac1{3a^2}\).
因此在 \(x=0\) 处取得极小值 \(0\)，在 \(x=\frac1a\) 处取得极大值 \(\frac1{3a^2}\).
\item 函数的正零点为 \(r=\frac{3}{2a}\)，且在 \(x>r\) 时 \(f(x)<0\).

若 \(0<a<\frac34\)，则 \(r>2\). 取 \(x_1\in(2,r)\) 并令其趋近于 \(r\)，则 \(f(x_1)>0\) 且可以任意接近 \(0\)，从而 \(1/f(x_1)\) 可以任意大. 但 \(x_2>1\) 时，\(f(x_2)\) 的正值不超过极大值 \(\frac1{3a^2}\)，不可能满足题设，故必须有 \(a\geqslant\frac34\).

当 \(\frac34\leqslant a\leqslant\frac32\) 时，\(1\leqslant r\leqslant2\). 对任意 \(x_1>2\)，都有 \(f(x_1)<0\). 在区间 \((r,+\infty)\subset(1,+\infty)\) 上，函数从 \(0\) 单调递减到 \(-\infty\)，值域为 \((-\infty,0)\)，所以总能取到负数 \(1/f(x_1)\)，题设成立.

若 \(a>\frac32\)，则 \(r<1\). 此时 \(x_2>1\) 时函数值域为 \((-\infty,f(1))\)，其中 \(f(1)=1-\frac{2a}{3}<0\). 而当 \(x_1\to+\infty\) 时，\(1/f(x_1)\to0^{-}\)，必有某些 \(x_1>2\) 使 \(1/f(x_1)>f(1)\)，该值不在 \(f((1,+\infty))\) 中，题设不成立. 因此 \(a\) 的取值范围为 \(\left[\frac34,\frac32\right]\).
\end{enumerate}
\end{solution}

\begin{problem}
已知 \(q\) 和 \(n\) 均为给定的大于 \(1\) 的自然数. 设集合 \(M = \{0, 1, 2, \cdots, q-1\}\)，集合 \(A = \{x \mid x = x_1 + x_2q + \cdots + x_nq^{n-1}, x_i \in M, i = 1, 2, \cdots, n\}\).

\begin{enumerate}
\item 当 \(q = 2\)，\(n = 3\) 时，用列举法表示集合 \(A\)；
\item 设 \(s\)，\(t \in A\)，\(s = a_1 + a_2q + \cdots + a_nq^{n-1}\)，\(t = b_1 + b_2q + \cdots + b_nq^{n-1}\)，其中 \(a_i\)，\(b_i \in M\)，\(i = 1\)，\(2\)，\(\cdots\)，\(n\). 证明：若 \(a_n < b_n\)，则 \(s < t\).
\end{enumerate}
\end{problem}

\begin{solution}
\begin{enumerate}
\item 当 \(q=2\)，\(n=3\) 时，每个 \(x_i\) 只能取 \(0\) 或 \(1\)，所以
\[
A=\{0,1,2,3,4,5,6,7\}
\]
\item 作差并利用 \(a_i\)，\(b_i\in\{0,1,\ldots,q-1\}\)，得
\[
t-s
=(b_n-a_n)q^{n-1}+\sum_{i=1}^{n-1}(b_i-a_i)q^{i-1}
\]
由于 \(b_n-a_n\geqslant1\)，且 \(b_i-a_i\geqslant-(q-1)\)，所以
\[
t-s\geqslant q^{n-1}-(q-1)\sum_{i=1}^{n-1}q^{i-1}=q^{n-1}-(q^{n-1}-1)=1>0
\]
因而 \(t-s>0\)，即 \(s<t\).
\end{enumerate}
\end{solution}
